\clearpage
\section{The question and the unit of evidence}\label{sec:question}
\subsection{Start with a concrete task}
Suppose a gallery contains one clean motion example for each of twenty actions.
A query is another motion clip. The system ranks the gallery by similarity and
returns the action of the closest example. A sitting-down query should retrieve
the sitting-down example even when some joint observations are unreliable.

This is \textbf{retrieval}: find the matching example. It is not a system that
estimates joint positions from raw video, predicts rowing force, or judges
whether an action was performed well. The input already contains estimated
joint positions. The action label provides the definition of a correct match;
it does not measure every possible meaning of biomechanical similarity.

\textbf{One-shot} describes the gallery: one labeled example for each new
action. It does not mean the model learns from one training clip. The comparison
rule is learned using thousands of clips from other actions, then applied to
actions excluded from this study's head training and model selection.

Write a motion example as $x$, its action label as $y$, and its learned
representation as $z=f(x)$. Let $\mathcal A$ be the twenty gallery action labels
and $g_a$ the reference clip for action $a$. Then
\[
 \widehat y(x)=\underset{a\in\mathcal A}{\arg\max}\;s(f(x),f(g_a)).
\]
The function $s$ is cosine similarity. Later sections derive both $f$ and $s$.
No final-test labels are used to update $f$.

\subsection{The hypothesis is deliberately narrower than ``better AI''}
There are three training recipes: standard contrastive, full contextual, and
supervised contrastive (SupCon). The one confirmatory comparison is contextual
versus standard contrastive. SupCon provides a descriptive reference.

Each trained head is evaluated on clean queries and on nine damaged-query
conditions. For a method $m$, seed $s$, and corruption condition $c$, define
\[
 d_{m,s,c}=A_{m,s,0}-A_{m,s,c},
\]
where $A$ is top-1 accuracy and $0$ denotes clean input. The main effect is
\[
 \Delta=\frac{1}{3\cdot9}\sum_{s,c}
       \bigl(d_{\mathrm{contextual},s,c}-d_{\mathrm{contrastive},s,c}\bigr).
\]
Negative $\Delta$ means the contextual recipe loses fewer accuracy points under
the declared damage mixture. A later section explains uncertainty and the
prespecified decision rule.

\begin{concept}{A smaller drop is not necessarily better damaged-input accuracy}
Invented example: standard contrastive goes from 90\% clean to 75\% damaged,
a 15-point drop. Contextual goes from 80\% to 70\%, a 10-point drop.
Then $\Delta=10-15=-5$ percentage points. Contextual degrades less, but its
damaged-input accuracy is five points worse. Both statements must be reported.
\end{concept}

The hypothesis concerns a \emph{full recipe}. Its contextual term, positive
margin, and similarity regularizer differ from the standard recipe. A positive
finding cannot identify the neighborhood term alone as the cause. The omitted
extra control limits causal attribution; it does not make the declared
recipe-to-recipe comparison meaningless.

\subsection{Why use this study design?}
Freezing an existing motion encoder makes the training comparison narrower and
less expensive than fine-tuning the whole network. Holding the head architecture,
data, duration, and tuning budget fixed makes differences easier to attribute to
the training recipes. Testing unfamiliar action classes asks whether the learned
comparison transfers beyond the action labels used for training.

Damaging only queries keeps the references unchanged. The experiment can then
measure sensitivity to the query disturbance. It does not establish performance
when both gallery and query are damaged, or robustness to arbitrary camera
failures. Those would be different experiments.

\source{\src{plan/research-plan.v3.md}{Final v3 plan, Sections 1, 3--6};
\src{docs/decisions/0004-frozen-one-shot-v3.md}{Decision 0004}.}

\section{Data: what one row actually means}\label{sec:data}
\subsection{Dataset, observations, and dependence}
The dataset family is NTU RGB+D 120,\cite{ntu} but this pipeline uses its adopted
\textbf{HRNet-derived 2D pose annotations}. The model channels are $x$ coordinate,
$y$ coordinate, and confidence. They are not $x,y,z$ ground-truth 3D positions.
``Clean'' means before this study adds damage; the pose estimates already have
their own errors.

The adopted inventory has 113,945 usable captures. Together with 535 documented
missing-skeleton exclusions, these reconcile to 114,480 nominal captures.
Reconciling the count matters because a changed input inventory could silently
change the population being compared.

A capture has metadata including setup, performer, repetition, action, and
camera. Several cameras can observe the \emph{same underlying performance}.
Thus three rows need not be three independent examples. The performance-group
key is
\[
  (\text{setup},\text{performer},\text{repetition},\text{action}),
\]
with camera deliberately omitted. This key matters twice: when excluding
near-duplicate gallery performances from primary queries, and when estimating
uncertainty.

\subsection{Split by action before learning the comparison}
The 120 action classes are partitioned into 100 auxiliary actions and twenty
novel actions. Within the auxiliary set, eighty actions train candidate heads
and twenty different actions support development evaluation.

\begin{center}
\resizebox{0.96\linewidth}{!}{%
\begin{tikzpicture}[node distance=8mm,
 box/.style={draw=ink,rounded corners=2pt,align=center,text width=37mm,
 minimum height=13mm,font=\small},arrow/.style={-{Latex},thick,ink}]
\node[box] (all) {120 action classes};
\node[box,below left=10mm and 10mm of all] (aux) {100 auxiliary actions};
\node[box,below right=10mm and 10mm of all] (novel) {20 novel actions\\final test};
\node[box,below left=10mm and -2mm of aux] (train) {80 actions\\candidate training};
\node[box,below right=10mm and -2mm of aux] (dev) {20 actions\\development retrieval};
\draw[arrow] (all)--(aux);\draw[arrow] (all)--(novel);
\draw[arrow] (aux)--(train);\draw[arrow] (aux)--(dev);
\end{tikzpicture}
}
\end{center}

\begin{tabularx}{\linewidth}{@{}Y r Y@{}}
\toprule
Role & Rows & Permitted use\\\midrule
Candidate-head training & 76,013 & Learn weights using eighty actions.\\
Development-validation pool & 18,988 & Twenty held-out auxiliary actions.\\
Development gallery & 20 & One fixed reference per development action.\\
Primary development queries & 18,929 & Evaluate candidates without gallery-performance views.\\
Final-head training & 95,001 & All 100 auxiliary actions, after selecting recipes.\\
Novel gallery & 20 & Fixed official exemplars.\\
Primary novel queries & 18,884 & Confirmatory final evaluation.\\
Exact-official novel queries & 18,924 & Descriptive compatibility evaluation.\\
\bottomrule
\end{tabularx}

These rows are \emph{roles}, not disjoint quantities to add together. Final-head
training is the union of the two auxiliary pools: $76{,}013+18{,}988=95{,}001$.
Gallery and query roles are selected from their respective validation or novel
pools. The twenty novel anchors plus 18,924 official queries give 18,944 novel
rows; $95{,}001+18{,}944=113{,}945$.

After choosing one learning rate for each recipe, nine fresh final heads are
trained on all 100 auxiliary actions. Reusing the development actions at this
point is legitimate because recipe selection is finished and the novel actions
still provide the final evaluation. The final heads are not the candidate or
pilot checkpoints with new names.

\subsection{Why two query definitions?}
The exact-official query set removes the gallery exemplar itself. The primary
set also removes other camera views of that exemplar's performance. Otherwise,
the system could retrieve a simultaneously recorded version of the same event,
which is easier than recognizing another performance of the action.

The two definitions answer related but different questions. The official one
maintains benchmark comparability; the primary one reduces this specific
cross-camera dependence. Their difference is not selected after seeing scores.
Exact IDs are locked from metadata. Do not infer counts by assuming every
performance has three valid camera views.

This is an \textbf{action-disjoint} study, not a claim of performer-disjoint
generalization. People or setups can recur. ``Novel'' describes the role in
this study's head training and selection; it does not prove the pretrained
encoder has never encountered any related action in any source data.

\subsection{A manifest is part of the experimental definition}
A manifest records which sample IDs belong to a role, their order, exclusions,
and source bindings. Its hash is a compact fingerprint of those exact bytes.
Recreating a manifest independently and comparing it checks that the stated
rules lead to the same population.

This is stronger than checking a row count. Two files can each contain 18,884
rows but describe different queries. Conversely, equal hashes establish equality
of content, not that the content is scientifically correct. Both semantic
validation and identity checks are needed.

\studycheck{Why would a random split of capture rows fail to test the intended
transfer to new actions? Why is removing only the exact gallery filename
insufficient? Answers: Questions 1--2 in the answer key.}
\source{\src{plan/research-plan.v3.md}{Final v3 plan, Section 3};
\src{docs/protocol/ntu-manifest-audit.v2.md}{Adopted manifest audit};
\src{docs/protocol/input-contract-amendment.v1.md}{Input contract amendment}.}
